\section{Analytical contrasts for bounded rationality}
\label{sec:annex-p17-rationality-8-5}

\begin{annexitem}{annex:p17.br.bias-variance}{thm:p17-br-bias-variance}{Independent test noise and the two specified risks}
\begin{proof}
Write \(\mu=E\widehat m\) and expand
\[
 \widehat m-Y=(\widehat m-\mu)+(\mu-m)-\epsilon.
\]
All terms have finite second moments. The expectation of
\(\widehat m-\mu\) and of \(\epsilon\) is zero. Independence
gives \(E[(\widehat m-\mu)\epsilon]=0\); the other mixed terms
also have zero expectation. Squaring and taking expectations therefore
gives the displayed decomposition. A constant predictor has zero variance
and squared bias \((a-m)^2\); an unbiased predictor has zero bias
and variance \(v\). Subtracting the common noise term proves the
strict comparison and gives a tie when the two quantities are equal.
For the stated values, the constant's risk is \(0^2+1=1\) when
\(m=0\) and \(3^2+1=10\) when \(m=3\); the unbiased
predictor's risk is \(4+1=5\) in both. A predictor taking values
\(m-2,m+2\) with equal probability realizes the stated mean and
variance. These examples require an independent future test response;
reusing the fitted response would not justify the vanishing mixed term.
\end{proof}
\end{annexitem}

\begin{annexitem}{annex:p17.br.noncompensation}{thm:p17-br-noncompensation}{A dominant binary cue and the tie boundary}
\begin{proof}
For an originally negative weight, write
\(w_ix_i=|w_i|(1-x_i)+w_i\). The constant term cancels between
the two alternatives, proving the stated orientation step. A
zero-weight coordinate contributes zero regardless of its value.

If no positive-weight cue discriminates, the weighted sum is zero.
Otherwise let \(r\) be the first discriminating positive-weight index.
Earlier terms vanish, and
\(d_r\in\{-1,1\}\). If \(r<k\),
\[
 \left|\sum_{j>r}w_jd_j\right|
          \leq\sum_{j>r}w_j<w_r=|w_rd_r|.
\]
The later terms cannot reverse or cancel the sign of \(w_rd_r\).
If \(r=k\), there are no later terms and \(w_r>0\) gives the
same conclusion. Thus the first discriminating cue and full linear
criterion have identical signs. An identical-vector tie must be resolved
consistently by the two procedures if a single selected alternative is
needed. Equality in the dominance bound is not sufficient for a strict
comparison: weights \((1,1)\) and differences \((1,-1)\) give a
zero full score despite a positive first cue. With weights \((1,2)\)
the same differences reverse the sign. The binary bounds, directions,
ordering and strict inequalities are therefore material assumptions.
The zero-weight convention is also necessary: with a single zero-weight
cue, differing attribute values still give a tied full score.
\end{proof}
\end{annexitem}

\begin{annexitem}{annex:p17.br.environment-contrast}{thm:p17-br-environment-contrast}{Resource accounting and changed environmental scores}
\begin{proof}
Cue~1 takes values one and zero, so the fixed procedure selects \(a\)
after one inspection. This meets the one-unit deadline. In \(N\),
the scores are \(4\cdot1+2\cdot0+1\cdot0=4\) and
\(4\cdot0+2\cdot1+1\cdot1=3\). In \(C\), they are
\(1\) and \(2+4=6\), so the same selected alternative loses.
Inspecting cue~3 instead shows zero for \(a\) and one for \(b\),
and selects \(b\) within the same one-unit budget. The explicitly
specified full-sum comparator acquires all three coordinates and uses
three time units, exceeding the deadline. These facts establish the
advertised feasibility and success/failure comparisons; they do not
claim that full summation is the only way of finding the right answer.
\end{proof}
\end{annexitem}
